%% kapa: exam=1 round=26 year=2015 period=1 form=A subject=civil_law subject_ko=민법 seq=8 number=8
%% source: Q-Net 감정평가사 기출문제 (cst00309) artlSeq=5138001 file=extracted/1008/5138001/제26회_감정평가사 1차_기출/제26회_감정평가사 1차_기출/1교시_감평_A형.pdf
%% conversion: kordoc 4.24.51 -> tools/gen1.py (auto); review=proofread-sonnet; flags=none
\begin{question}{8}
의사표시에 관한 설명으로 옳지 않은 것은? (다툼이 있으면 판례에 따름)
\begin{choices}
  \item 의사표시자가 그 통지를 발송한 후 사망하여도 의사표시의 효력에 영향을 미치지 않는다.
  \item 통정허위표시의 무효로 대항할 수 없는 제3자는 허위표시의 당사자와 그의 포괄승계인 이외의 자로서 허위표시행위를 기초로 하여 새로운 이해관계를 맺은 자를 말한다.
  \item 진의 아닌 의사표시에 있어서의 ‘진의’란 특정한 내용의 의사표시를 하고자 하는 표의자의 생각을 말하는 것이지 표의자가 진정으로 마음 속에서 바라는 사항을 뜻하는 것은 아니다.
  \item 통정한 허위표시에 의하여 외형상 형성된 법률관계로 생긴 채권을 가압류한 자는 통정허위표시의 무효로 대항할 수 없는 제3자에 해당한다.
  \item 제3자가 통정허위표시의 무효에 대항하기 위해서는 선의ㆍ무과실이어야 한다.
\end{choices}
\end{question}
\answer{5}
